\section{Concept Hypothesis}
\label{sec:concept_hypothesis}

In this section, we introduce the main working hypothesis underlying much of the literature on \gls{ls} analysis.  
Although rarely stated explicitly, this hypothesis is implicitly present in the recurring use of notions such as abstraction, representation, \gls{ls}, concept, and semantics.
Throughout this manuscript, we refer to this hypothesis as \ConceptHypothesisName.
The goal of this section is to make this conceptual framework explicit by defining its main components.  

The \ConceptHypothesisName originates from the need to explain the generalization ability of deep \gls{nn} observed after training, as discussed in Section~\ref{sec:evaluation}.
A model is said to generalize when, for an input $x \in \mathcal{X}$ not seen during training, it still produces an output $\hat{y}$ close to the true output $y$.  
Such behavior cannot be explained by simple memorization of training samples.  
Instead, the model must extract and exploit underlying regularities in the target function.  
This requires a process of abstraction, as defined in Definition~\ref{def:abstraction}, whereby raw inputs are transformed into internal structures that are meaningful for the task \cite{bengio2014representationlearningreviewnew}.  

\begin{definition}{Abstraction}{abstraction}
Abstraction is the process by which a model transforms an input into a task-relevant internal form.  
It captures and restructures information in order to represent relationships that are useful for prediction.  
\end{definition}

The ability of \gls{nn} to perform abstraction arises from their intermediate activations.
A \gls{nn} does not directly map $x$ to $\hat{y}$, it computes a sequence of intermediate activations across hidden layers.
These activations are numerical encodings of the input referred to as embeddings.
Once training is complete, these embeddings are structured to support the performance of a specific task.
To mark this distinction, we will refer to these trained embeddings as \gls{lr}, as defined in Definition~\ref{def:representation}.

\begin{definition}{Latent Representation}{representation}
A \gls{lr} is the internal encoding of an input given by an intermediate layer of a trained \gls{nn}.
\end{definition}

Each hidden layer refines the \gls{lr} of the previous layer.
Each transformation progressively reshapes the \gls{lr} into a form more directly aligned with the output space \cite{rumelhart1986learning}.
This process can be described as a sequence of transformations where each $\latentRepresentation[l]$ is a \gls{lr}:
\[
x \rightarrow \latentRepresentation[1] \rightarrow \latentRepresentation[2] \rightarrow \dots \rightarrow \latentRepresentation[L-1] \rightarrow \hat{y}
\]

The set of all such \gls{lr} produced by a layer over a dataset defines its \gls{ls}, as defined in Definition~\ref{def:latent_space}.

\begin{definition}{Latent Space}{latent_space}
The \gls{ls} is the set of all \gls{lr} produced by a layer over a dataset.
It is shaped by learning and organizes \gls{lr} according to task-relevant properties.
\end{definition}

Under the \ConceptHypothesisName, the analysis of the \gls{ls} is assumed to provide access to abstract structures learned by the model, as illustrated in Figure~\ref{fig:latent_space}.
The same set of \gls{lr} can be viewed either as an unstructured point cloud or, under this hypothesis, as a collection of points organized by concept membership, in which case the structure of the \gls{ls} becomes meaningful and amenable to analysis.
In Figure~\ref{fig:latent_space}, the colors represent hypothetical concept memberships assigned by the model.
In this manuscript, a concept is defined as an abstract representation that organizes information and supports reasoning (see Definition~\ref{def:concept}).

\documentclass[tikz, border=10pt]{standalone}
\usepackage{tikz}
\usepackage{pgfplots}
\usepackage{pgfkeys}
\usetikzlibrary{arrows.meta}
\pgfplotsset{compat=1.18}

% ================================================================
%  \plotPoints — Reusable macro to render a point cloud from file
%
%  Usage:
%    \plotPoints[color=<c>, size=<s>, opacity=<o>]{<file.data>}
%
%  Arguments (all optional, defaults shown below):
%    color   : fill & stroke color of the markers  (default: black)
%    size    : marker radius                        (default: 2pt)
%    opacity : fill opacity [0..1]                  (default: 0.85)
%
%  Positional (mandatory):
%    #1 : path to a whitespace-separated x y data file
%         (lines starting with % are treated as comments)
%
%  Example:
%    \plotPoints[color=blue!70!cyan, size=3pt]{data/cluster_a.data}
% ================================================================
\pgfkeys{
  /plotPoints/.is family, /plotPoints,
  % --- defaults ---
  color/.initial   = black,
  size/.initial    = 2pt,
  opacity/.initial = 0.85,
}

\newcommand{\plotPoints}[2][]{%
  \pgfkeys{/plotPoints, #1}%                     % apply user overrides
  \pgfkeysgetvalue{/plotPoints/color}  {\ppColor}%
  \pgfkeysgetvalue{/plotPoints/size}   {\ppSize}%
  \pgfkeysgetvalue{/plotPoints/opacity}{\ppOpacity}%
  \addplot[
    only marks,
    mark        = *,
    mark size   = \ppSize,
    mark options= {fill=\ppColor, draw=\ppColor, fill opacity=\ppOpacity},
  ] file {#2};%
}
% ================================================================
%  \plotEllipse — Companion macro to draw a confidence ellipse
%
%  Usage:
%    \plotEllipse[color=<c>, width=<w>]
%                {<cx>}{<cy>}{<angle>}{<rx>}{<ry>}
%
%  Arguments (all optional):
%    color : stroke color  (default: black)
%    width : line width    (default: very thick)
%
%  Positional (mandatory):
%    {cx}    : ellipse center x  (axis coordinates)
%    {cy}    : ellipse center y  (axis coordinates)
%    {angle} : counter-clockwise rotation in degrees
%    {rx}    : horizontal semi-axis (cm)
%    {ry}    : vertical   semi-axis (cm)
%
%  Example:
%    \plotEllipse[color=orange]{-3}{1.5}{15}{2.3}{1.2}
% ================================================================
\pgfkeys{
  /plotEllipse/.is family, /plotEllipse,
  color/.initial = black,
  width/.initial = very thick,
}

\newcommand{\plotEllipse}[6][]{%
  \pgfkeys{/plotEllipse, #1}%
  \pgfkeysgetvalue{/plotEllipse/color}{\peColor}%
  \pgfkeysgetvalue{/plotEllipse/width}{\peWidth}%
  \draw[\peColor, \peWidth, rotate around={#4:(#2,#3)}]
    (#2,#3) ellipse (#5cm and #6cm);%
}

% ================================================================
%  Document
% ================================================================
\begin{document}
\begin{tikzpicture}

\begin{axis}[
  width        = 11cm,
  height       = 10cm,
  xmin=-7, xmax=8,
  ymin=-7, ymax=5.5,
  % --- Arrow-only axes: no box, no labels, no ticks ---
  axis x line  = bottom,
  axis y line  = left,
  axis line style = {
    -{Stealth[length=8pt, width=6pt]},
    thick,
  },
  ticks        = none,
  xlabel       = {},
  ylabel       = {},
  clip         = true,
]

% ----------------------------------------------------------------
%  Cluster A
%  Generator : numpy.random.multivariate_normal, seed=42
%  Mean      : mu    = (-3.0,  1.5)
%  Covariance: sigma = (1.8, 0.9), angle = 15 deg
%              => R * diag(1.8^2, 0.9^2) * R^T
%  N points  : 45
%  Ellipse   : rotation=15 deg, semi-axes=(2.3cm, 1.2cm)
% ----------------------------------------------------------------
% \plotEllipse[color=orange]{-3}{1.5}{15}{2.3}{1.2}
\plotPoints{cluster_a.data}

% ----------------------------------------------------------------
%  Cluster B
%  Generator : numpy.random.multivariate_normal, seed=42
%  Mean      : mu    = (-2.5, -4.0)
%  Covariance: sigma = (1.7, 0.85), angle = 10 deg
%              => R * diag(1.7^2, 0.85^2) * R^T
%  N points  : 40
%  Ellipse   : rotation=10 deg, semi-axes=(2.2cm, 1.1cm)
% ----------------------------------------------------------------
% \plotEllipse[color=green!50!black]{-2.5}{-4}{10}{2.2}{1.1}
\plotPoints{cluster_b.data}

% ----------------------------------------------------------------
%  Cluster C
%  Generator : numpy.random.multivariate_normal, seed=42
%  Mean      : mu    = (4.2, -0.5)
%  Covariance: sigma = (1.1, 2.0), angle = -20 deg
%              => R * diag(1.1^2, 2.0^2) * R^T
%  N points  : 38  (2 outliers beyond axis bounds removed)
%  Ellipse   : rotation=-20 deg, semi-axes=(1.4cm, 2.6cm)
% ----------------------------------------------------------------
% \plotEllipse[color=violet]{4.2}{-0.5}{-20}{1.4}{2.6}
\plotPoints{cluster_c.data}

\end{axis}
\end{tikzpicture}
\end{document}

\begin{definition}{Concept}{concept}
A concept is a task-dependent abstraction that captures regularities in the input space relevant for predicting the target output.
\end{definition} 

Our definition of concepts is consistent with both the pragmatic tradition in philosophy and the information-processing view of cognition.
Within these frameworks, concepts arise from the need to simplify an otherwise intractable reality.
Because the environment contains more information than any cognitive system can process, intelligent behavior relies on abstractions, called concepts, that discard irrelevant details while preserving the information required to achieve a given objective \cite{Peirce1878,Misak2013}.
These concepts are realized as representations that reduce complexity and enable reasoning under limited computational resources \cite{Newell1980,Simon1996}.
Although these theories were originally developed to explain human cognition, the \ConceptHypothesisName assumes that the same principles can be extended to deep \gls{nn}.

In order to influence the model's computations, concepts must necessarily admit a concrete realization within the \gls{ls}.
We refer to this realization as the \gls{cs}.
The definition of the \gls{cs} is provided in Definition~\ref{def:concept_structure}.

\begin{definition}{Concept Structure}{concept_structure}
A \gls{cs} is the specific form for how a concept is instantiated within the \gls{ls}.
\end{definition}

The methods reviewed in this chapter all build on the \ConceptHypothesisName formalized in Hypothesis~\ref{hyp:concept}, but differ in how they characterize the \gls{cs} of a concept.
This choice directly determines how concepts can be identified and extracted from the \gls{ls}.
As a result, several families of \gls{ls} analysis methods have emerged in the literature, each adopting a different view of how concepts are structured within the \gls{ls}.

\begin{hypothesis}{Concept Hypothesis}{concept}
\gls{nn} generalize by performing abstraction.
Through training, their \gls{ls} becomes organized around concepts.
Concepts are task-dependent abstractions that capture regularities relevant to the task.
Analyzing the \gls{ls} therefore provides access to the \gls{cs} through which these concepts are instantiated in the model.
\end{hypothesis}

Now that the definition of a \ConceptHypothesisName has been established, we can consider its relationship to human.
According to our definition, we can imagine concepts that are useful for the machine but do not make sense to humans.
This relationship is captured by the notion of semantics \cite{marconato2023interpretabilitymindbeholdercausal}, defined in Definition~\ref{def:semantics}.
Some concepts developed in \gls{ls} can be readily associated with human concepts and are said to have high semantic value.
Others may contribute to task performance while remaining difficult for humans to understand, and are said to have low semantic value.

\begin{definition}{Semantics}{semantics}
Semantics refers to the alignment between concepts in the model \gls{ls} and concepts used by a human observer.  
A concept is semantic if it can be associated with a human interpretable concept.  
\end{definition}

It is important to note that this conceptual framework is not supported by a complete theoretical understanding of deep \gls{nn}.
Current mathematical tools do not provide a rigorous characterization of their internal representations and do not formally prove that concepts emerge as identifiable \gls{cs} within \gls{ls}.
Nevertheless, this perspective constitutes a widely adopted working hypothesis in the field of representation and \gls{ls} analysis.

The remainder of this chapter is devoted to reviewing methods for analyzing \gls{ls} and extracting the concepts they encode.
We begin with probing, one of the most direct approaches for assessing whether concepts are represented in the \gls{ls}.